\label{kap:financni}
\mathquote{Compound interest is the eighth wonder of the world. He who understands it, earns it; he who doesn't, pays it.}{připisováno Albertu Einsteinovi (nejspíš neprávem)}

Finanční matematika je ta část matematiky, která vás bude v životě pronásledovat nejdéle. Spořicí účty, hypotéky, splátky telefonu, kreditní karty, inflace. Banky a obchodníci s tím počítají každý den a spoléhají na to, že vy ne. Matematicky v této kapitole nic nového nebude: všechno jsou to procenta, geometrické posloupnosti z kapitoly~\ref{kap:posloupnosti} a exponenciála z kapitoly~\ref{kap:exponencialni}. Nové bude jen to, že tentokrát jde o vaše peníze.

\section{Procenta}\label{sec:procenta}

\begin{definitionbox}
  \textbf{Procento} je setina celku: $1\,\% = \frac{1}{100} = 0{,}01$. \textbf{Promile} je tisícina: $0{,}1\,\%$, tedy $0{,}001$.

  \vspace{0.2cm}

  \begin{itemize}
    \item $p\,\%$ ze základu $z$ je \textbf{procentová část} $c = \frac{p}{100}\cdot z$,
    \item zvětšení o $p\,\%$ znamená vynásobení číslem $1 + \frac{p}{100}$, zmenšení o $p\,\%$ vynásobení číslem $1 - \frac{p}{100}$,
    \item \textbf{procentní bod} je rozdíl dvou procentních údajů. Pokud úroková sazba vzroste ze $2\,\%$ na $3\,\%$, vzrostla o $1$ procentní bod, ale o $50\,\%$.
  \end{itemize}
\end{definitionbox}

\begin{example}
  Bunda stála $2\,000$ Kč. Obchod ji zdražil o $20\,\%$ a pak ji ve výprodeji zlevnil o $20\,\%$. Kolik stojí teď?
  \begin{equation*}
    2000\cdot 1{,}2\cdot 0{,}8 = 2000\cdot 0{,}96 = 1920\ \text{Kč}.
  \end{equation*}
  Bunda je o $4\,\%$ levnější než na začátku, ne stejně drahá. Zlevnění se počítá z vyšší ceny, takže ubere víc, než předtím přidalo zdražení. Obchody to vědí. Zákazníci většinou ne.
  \begin{center}
    \begin{tikzpicture}[xscale=0.9, yscale=0.6]
      \draw[->] (0,0) -- (7.2,0);
      \draw[->] (0,0) -- (0,5.3) node[above] {\small Kč};
      \foreach \y/\l in {1/500, 2/1000, 3/1500, 4/2000} \draw (0.07,\y) -- (-0.07,\y) node[left] {\tiny $\l$};
      \filldraw[fill=gray!35, draw=black] (0.6,0) rectangle (1.8,4);
      \filldraw[fill=gray!35, draw=black] (2.9,0) rectangle (4.1,4);
      \filldraw[fill=colorofwisdom!60, draw=black] (2.9,4) rectangle (4.1,4.8);
      \filldraw[fill=gray!35, draw=black] (5.2,0) rectangle (6.4,3.84);
      \draw[dashed, colorofwisdom] (5.2,3.84) rectangle (6.4,4.8);
      \draw[dashed, gray] (0,4) -- (7,4);
      \node[font=\scriptsize] at (3.5,4.4) {$+400$};
      \node[font=\scriptsize, colorofwisdom] at (5.8,4.32) {$-480$};
      \node[below, font=\small] at (1.2,0) {původně};
      \node[below, font=\small] at (3.5,0) {$+20\,\%$};
      \node[below, font=\small] at (5.8,0) {$-20\,\%$};
      \node[above, font=\scriptsize] at (1.2,4) {$2000$};
      \node[above, font=\scriptsize] at (3.5,4.8) {$2400$};
      \node[below, font=\scriptsize] at (5.8,3.84) {$1920$};
    \end{tikzpicture}
  \end{center}
  Na obrázku je to vidět: oba kroky jsou `$20\,\%$', ale pokaždé z jiného základu. Zdražení přidá $20\,\%$ ze $2\,000$, tedy $400$ Kč, zlevnění ubere $20\,\%$ ze $2\,400$, tedy $480$ Kč.
\end{example}

\begin{example}
  Po zdražení o $25\,\%$ stojí lístek do kina $250$ Kč. Kolik stál předtím? Ne $250\cdot 0{,}75 = 187{,}5$ Kč, ale $\frac{250}{1{,}25} = 200$ Kč. Abyste zrušili zdražení o $25\,\%$, musíte zlevnit o $20\,\%$.
\end{example}

\begin{problem}
  \begin{enumerate}[label=(\alph*)]
    \item Cena akcie vzrostla o $10\,\%$, pak klesla o $10\,\%$, pak vzrostla o $10\,\%$. Jaká je celková změna?
    \item Kolika procenty je $45$ z $180$? A kolik je $180$ procent z $45$?
    \item Volební strana měla $8\,\%$ a teď má $10\,\%$. O kolik procentních bodů a o kolik procent si polepšila?
    \item O kolik procent se zvětší obsah čtverce, když jeho stranu zvětšíme o $10\,\%$?
  \end{enumerate}
\end{problem}

\begin{solution}
  \begin{enumerate}[label=(\alph*)]
    \item $1{,}1\cdot 0{,}9\cdot 1{,}1 = 1{,}089$, tedy vzrostla o $8{,}9\,\%$.
    \item $\frac{45}{180} = 0{,}25 = 25\,\%$. $1{,}8\cdot 45 = 81$.
    \item O $2$ procentní body, což je $\frac{2}{8} = 25\,\%$.
    \item $1{,}1^2 = 1{,}21$, tedy o $21\,\%$.
  \end{enumerate}
\end{solution}

\section{Úročení}\label{sec:uroceni}

\begin{definitionbox}
  Pojmy:
  \begin{itemize}
    \item \textbf{jistina} $K_0$: počáteční vklad (nebo dluh),
    \item \textbf{úroková sazba} $i$: kolik procent z jistiny se připíše za jedno úrokovací období (obvykle rok, zkratka p.\,a., \textit{per annum}). Ve vzorcích ji píšeme jako desetinné číslo, tedy $3\,\% = 0{,}03$,
    \item \textbf{úrok}: částka, kterou za půjčení peněz dostaneme (nebo zaplatíme).
  \end{itemize}

  \textbf{Jednoduché úročení}: úrok se počítá stále jen z původní jistiny. Po $n$ obdobích
  \begin{equation*}
    K_n = K_0(1 + n\cdot i).
  \end{equation*}

  \textbf{Složené úročení}: úrok se připíše k jistině a v dalším období se úročí i on. Po $n$ obdobích
  \begin{equation*}
    K_n = K_0{(1 + i)}^n.
  \end{equation*}
\end{definitionbox}

Jednoduché úročení je aritmetická posloupnost (každý rok přibude stejná částka), složené úročení je geometrická posloupnost s kvocientem $1 + i$ (každý rok se částka vynásobí stejným číslem). Na krátkou dobu je rozdíl malý, na dlouhou obrovský, protože geometrická posloupnost roste exponenciálně.

Představte si to takhle. Při jednoduchém úročení dostáváte každý rok stejnou `výplatu' z původní jistiny. Při složeném úročení pracují i úroky z minulých let: úrok z úroku, pak úrok z úroku z úroku a tak dál. Na grafu je vývoj vkladu $10\,000$ Kč při $3\,\%$ ročně z následujícího příkladu, vždy po pěti letech. Šedé sloupce rostou rovnoměrně (přímka), červené stále rychleji (exponenciála).

\begin{center}
  \begin{tikzpicture}[yscale=1.1, xscale=0.95]
    \draw[->] (0,0) -- (11.3,0) node[below right] {\small roky};
    \draw[->] (0,0) -- (0,3.9) node[above] {\small tis. Kč};
    \foreach \y in {1,2,3} {
      \draw (0.07,\y) -- (-0.07,\y) node[left] {\tiny $\y0$};
      \draw[gray!30] (0.07,\y) -- (11,\y);
    }
    \foreach \k in {0,...,8} {
      \pgfmathsetmacro{\n}{5*\k}
      \pgfmathsetmacro{\s}{1 + 0.03*\n}
      \pgfmathsetmacro{\c}{pow(1.03,\n)}
      \filldraw[fill=gray!40, draw=black] ({1.2*\k+0.25},0) rectangle ({1.2*\k+0.6},\s);
      \filldraw[fill=colorofwisdom!70, draw=black] ({1.2*\k+0.6},0) rectangle ({1.2*\k+0.95},\c);
      \pgfmathtruncatemacro{\nn}{\n}
      \node[below, font=\tiny] at ({1.2*\k+0.6},0) {$\nn$};
    }
    \node[right, font=\scriptsize] at (10.6,2.2) {$22\,000$};
    \node[right, font=\scriptsize, colorofwisdom] at (10.6,3.26) {$32\,620$};
    \filldraw[fill=gray!40, draw=black] (0.5,3.55) rectangle (0.8,3.75);
    \node[right, font=\scriptsize] at (0.8,3.65) {jednoduché};
    \filldraw[fill=colorofwisdom!70, draw=black] (3.0,3.55) rectangle (3.3,3.75);
    \node[right, font=\scriptsize] at (3.3,3.65) {složené};
  \end{tikzpicture}
\end{center}

\begin{example}
  Vložíme $10\,000$ Kč na $4$ roky s úrokem $3\,\%$ ročně.
  \begin{itemize}
    \item Jednoduché úročení: $10000(1 + 4\cdot 0{,}03) = 11\,200$ Kč.
    \item Složené úročení: $10000\cdot 1{,}03^4\approx 11\,255{,}09$ Kč.
  \end{itemize}
  Za čtyři roky je rozdíl jen $55$ Kč. Za čtyřicet let by to bylo $22\,000$ Kč proti $32\,620$ Kč.
\end{example}

\begin{definitionbox}
  \textbf{Daň z úroků}: V České republice se z úroků na běžných účtech odvádí \textbf{srážková daň} $15\,\%$. Ze sazby $i$ nám tedy reálně zbude $0{,}85\cdot i$ a
  \begin{equation*}
    K_n = K_0{(1 + 0{,}85\cdot i)}^n.
  \end{equation*}

  \textbf{Častější připisování úroků}: Pokud je roční (nominální) sazba $i$ a úroky se připisují $m$-krát ročně, připíše se pokaždé $\frac{i}{m}$ a za $n$ let
  \begin{equation*}
    K_n = K_0{\left(1 + \frac{i}{m}\right)}^{mn}.
  \end{equation*}
  Skutečný roční výnos (\textbf{efektivní sazba}) je pak ${\left(1 + \frac{i}{m}\right)}^m - 1$, což je víc než $i$.
\end{definitionbox}

\begin{example}\label{ex:efektivni-sazba}
  Při sazbě $3\,\%$ a dani $15\,\%$ máme z $10\,000$ Kč za $4$ roky $10000\cdot 1{,}0255^4\approx 11\,059{,}68$ Kč.

  Při nominální sazbě $6\,\%$ s měsíčním připisováním je efektivní sazba ${(1 + 0{,}005)}^{12} - 1\approx 6{,}17\,\%$. Pokud bychom úroky připisovali nekonečně často (spojité úročení), dostali bychom efektivní sazbu $e^{0{,}06} - 1\approx 6{,}18\,\%$. Proč právě $e$? Připomeňme definici z kapitol~\ref{kap:exponencialni} a~\ref{kap:posloupnosti}:
  \begin{equation*}
    e = \lim_{n\to\infty}{\left(1 + \frac{1}{n}\right)}^n.
  \end{equation*}
  Pokud v našem výrazu označíme $m = n\cdot i$ (tedy počet připsání za rok je $n$-krát roční sazba), dostaneme
  \begin{equation*}
    {\left(1 + \frac{i}{m}\right)}^m = {\left(1 + \frac{1}{n}\right)}^{n\cdot i} = {\left({\left(1 + \frac{1}{n}\right)}^n\right)}^i\ \longrightarrow\ e^i\qquad\text{pro } m\to\infty.
  \end{equation*}
  Při spojitém úročení tedy jistina roste přesně jako exponenciála: $K(t) = K_0e^{it}$. Tady se Eulerovo číslo z kapitoly~\ref{kap:exponencialni} objevuje přesně v tom kontextu, ve kterém ho objevil Jacob Bernoulli.
\end{example}

\begin{note}
  \textbf{Pravidlo 72}: Doba (v letech), za kterou se vklad při složeném úročení zdvojnásobí, je přibližně $\frac{72}{p}$, kde $p$ je sazba v procentech. Při $7\,\%$ je to asi $10$ let (přesně $\frac{\ln 2}{\ln 1{,}07}\approx 10{,}24$). Platí to proto, že $\ln 2\approx 0{,}693$ a pro malá $i$ je $\ln(1 + i)\approx i$. Číslo $72$ se používá místo $69$ proto, že se lépe dělí (dvěma, třemi, čtyřmi, šesti, osmi, devíti a dvanácti).


  Pravidlo skrývá ještě jednu pěknou věc: doba zdvojnásobení \textit{nezávisí na tom, kolik zrovna máte}. Z $1$ na $2$ to trvá stejně dlouho jako ze $2$ na $4$ nebo ze $4$ na $8$. Exponenciála je jediná křivka s touto vlastností, a proto se po stejných časových krocích stále zdvojnásobuje:
  \begin{center}
    \begin{tikzpicture}[xscale=0.3, yscale=0.42]
      \draw[->] (0,0) -- (34,0) node[right] {\small roky};
      \draw[->] (0,0) -- (0,9.4) node[above] {\small $K_n/K_0$};
      \foreach \y in {1,2,4,8} {
        \draw[dashed, gray] (0,\y) -- ({ln(\y)/ln(1.07)},\y) -- ({ln(\y)/ln(1.07)},0);
        \node[left, font=\scriptsize] at (0,\y) {$\y$};
      }
      \foreach \t in {10,20,30} \node[below, font=\scriptsize] at (\t,0) {$\t$};
      \draw[thick, colorofwisdom, domain=0:32.5, samples=60] plot (\x, {pow(1.07,\x)});
      \foreach \a/\b in {0/10.24, 10.24/20.49, 20.49/30.73} {
        \draw[<->, colorofwisdom] (\a,-1.6) -- (\b,-1.6);
        \node[below, font=\tiny] at ({(\a+\b)/2},-1.6) {$\approx 10$ let};
      }
      \node[font=\small] at (8,7) {$1{,}07^n$};
    \end{tikzpicture}
  \end{center}
\end{note}

\begin{definitionbox}
  \textbf{Inflace} je růst cenové hladiny. Pokud je úroková sazba $i$ a inflace $\pi$, je \textbf{reálný výnos}
  \begin{equation*}
    r = \frac{1 + i}{1 + \pi} - 1\approx i - \pi.
  \end{equation*}
\end{definitionbox}

\begin{example}
  Banka nabízí $4\,\%$, inflace je $6\,\%$. Reálný výnos je $\frac{1{,}04}{1{,}06} - 1\approx -1{,}89\,\%$. Peníze na účtu se sice množí, ale koupíte si za ně každý rok méně. A to jsme ještě nezapočítali daň.
\end{example}

\begin{problem}
  \begin{enumerate}[label=(\alph*)]
    \item Jakou roční sazbu potřebujeme, aby se vklad za $10$ let zdvojnásobil?
    \item Auto za $500\,000$ Kč ztrácí každý rok $15\,\%$ své hodnoty. Kolik bude stát za $4$ roky?
    \item Kolik musíme dnes vložit při sazbě $5\,\%$ (bez daně), abychom za $18$ let měli $100\,000$ Kč?
    \item Babička vám v den vašeho narození uložila $10\,000$ Kč na $4\,\%$ ročně (bez daně). Kolik budete mít k maturitě (v $19$ letech)?
  \end{enumerate}
\end{problem}

\begin{solution}
  \begin{enumerate}[label=(\alph*)]
    \item ${(1 + i)}^{10} = 2$, $i = \sqrt[10]{2} - 1\approx 7{,}18\,\%$.
    \item $500000\cdot 0{,}85^4\approx 261\,003$ Kč. Auto ztratí skoro polovinu hodnoty.
    \item $K_0 = \frac{100000}{1{,}05^{18}}\approx 41\,552$ Kč. Tomuto se říká \textbf{současná hodnota} budoucí částky.
    \item $10000\cdot 1{,}04^{19}\approx 21\,068$ Kč. Za tu dobu se vklad zhruba zdvojnásobí, přesně podle pravidla 72 ($\frac{72}{4} = 18$).
  \end{enumerate}
\end{solution}

\section{Spoření}\label{sec:sporeni}

Při pravidelném spoření vkládáme každé období stejnou částku $a$. Každý vklad se pak úročí jinou dobu: ten první nejdéle, ten poslední nejkratší dobu nebo vůbec. Výsledná částka je součet geometrické posloupnosti.

\begin{center}
  \begin{tikzpicture}[scale=1]
    \draw[->] (-0.3,0) -- (10.6,0) node[right] {\small čas};
    \foreach \k in {0,...,5} {
      \draw (2*\k,0.1) -- (2*\k,-0.1) node[below] {\scriptsize $\k$};
    }
    \foreach \k/\e in {1/4, 2/3, 3/2, 4/1} {
      \filldraw[colorofwisdom] (2*\k,0) circle (2pt);
      \node[above, font=\scriptsize] at (2*\k,0.05) {$a$};
      \draw[->, colorofwisdom!80, looseness=0.9] (2*\k,0.35) to[out=60, in=180] (9.6,{0.35+0.45*\e});
      \ifnum\e=1
        \node[right, font=\scriptsize] at (10,{0.35+0.45*\e}) {$a(1 + i)$};
      \else
        \node[right, font=\scriptsize] at (10,{0.35+0.45*\e}) {$a{(1 + i)}^{\e}$};
      \fi
    }
    \filldraw[colorofwisdom] (10,0) circle (2pt);
    \node[right, font=\scriptsize] at (10,0.35) {$a$};
    \node[font=\scriptsize, align=center] at (0,-0.9) {začátek};
    \node[font=\scriptsize, align=center] at (10,-0.9) {konec};
  \end{tikzpicture}
\end{center}
Na obrázku je spoření po dobu $n = 5$ období. Každý vklad `jede' do konce jinak dlouho: první vklad se úročí čtyři období, poslední vůbec. Na konci sečteme pět čísel $a, a(1 + i), \ldots, a{(1 + i)}^4$, a to je geometrická posloupnost.

\begin{definitionbox}
  \textbf{Pravidelné spoření}: Pokud na konci každého z $n$ období vložíme částku $a$ a sazba za období je $i$, máme na konci
  \begin{equation*}
    S_n = a\cdot\frac{{(1 + i)}^n - 1}{i}.
  \end{equation*}
  Pokud vkládáme na \textbf{začátku} každého období, každý vklad se úročí o jedno období déle a
  \begin{equation*}
    S_n = a{(1 + i)}\cdot\frac{{(1 + i)}^n - 1}{i}.
  \end{equation*}
\end{definitionbox}

\begin{proof}
  Při vkladech na konci období se poslední vklad neúročí vůbec, předposlední jedno období a první $n - 1$ období. Celkem
  \begin{equation*}
    S_n = a + a(1 + i) + a{(1 + i)}^2 + \cdots + a{(1 + i)}^{n-1}.
  \end{equation*}
  To je součet $n$ členů geometrické posloupnosti s prvním členem $a$ a kvocientem $q = 1 + i$. Podle kapitoly~\ref{kap:posloupnosti}
  \begin{equation*}
    S_n = a\cdot\frac{{(1 + i)}^n - 1}{(1 + i) - 1} = a\cdot\frac{{(1 + i)}^n - 1}{i}.
  \end{equation*}
\end{proof}

\begin{example}
  Každý měsíc na konci měsíce vložíme $1\,000$ Kč na účet s roční sazbou $3\,\%$ (měsíčně tedy $0{,}25\,\%$, daň zanedbáme). Kolik budeme mít za dva roky?
  \begin{equation*}
    S_{24} = 1000\cdot\frac{1{,}0025^{24} - 1}{0{,}0025}\approx 24\,703\ \text{Kč}.
  \end{equation*}
  Vložili jsme $24\,000$ Kč, úroky dělají $703$ Kč. Na krátkou dobu se spoření moc nevyplatí. Na čtyřicet let by to ale bylo přes $925\,000$ Kč z vložených $480\,000$ Kč.
\end{example}

\section{Úvěry}\label{sec:uvery}

Úvěr je spoření obráceně. Banka nám půjčí částku $U$ a my ji splácíme pravidelnými splátkami. Každá splátka se skládá z \textbf{úroku} (z dosud nesplaceného dluhu) a \textbf{úmoru} (části, která skutečně snižuje dluh).

\begin{definitionbox}
  \textbf{Anuitní splácení}: Úvěr $U$ splácíme $n$ stejnými splátkami $a$ na konci každého období při sazbě $i$ za období. Pak
  \begin{equation*}
    a = U\cdot\frac{i{(1 + i)}^n}{{(1 + i)}^n - 1}.
  \end{equation*}
\end{definitionbox}

\begin{proof}
  Dluh po jednom období je $U(1 + i) - a$, po dvou $\left(U(1 + i) - a\right)(1 + i) - a = U{(1 + i)}^2 - a(1 + i) - a$, a tak dál. Po $n$ obdobích má být dluh nulový:
  \begin{equation*}
    U{(1 + i)}^n - a\left({(1 + i)}^{n-1} + \cdots + (1 + i) + 1\right) = 0.
  \end{equation*}
  Závorka je stejná geometrická řada jako u spoření, tedy $\frac{{(1 + i)}^n - 1}{i}$. Odtud
  \begin{equation*}
    U{(1 + i)}^n = a\cdot\frac{{(1 + i)}^n - 1}{i}\quad\Leftrightarrow\quad a = U\cdot\frac{i{(1 + i)}^n}{{(1 + i)}^n - 1}.
  \end{equation*}
  Rovnice má mimochodem pěknou interpretaci: to, na kolik by narostl dluh, kdybychom nesplácali (vlevo), se musí rovnat tomu, na kolik by narostly splátky, kdybychom je místo toho spořili (vpravo).
\end{proof}

\begin{example}\label{ex:splatkovy-kalendar}
  Půjčíme si $200\,000$ Kč na $5$ let s úrokem $6\,\%$ ročně a splácíme jednou ročně. Roční splátka je
  \begin{equation*}
    a = 200000\cdot\frac{0{,}06\cdot 1{,}06^5}{1{,}06^5 - 1}\approx 47\,479{,}28\ \text{Kč}.
  \end{equation*}
  Celkem zaplatíme $5a\approx 237\,396$ Kč, z toho úroky jsou asi $37\,396$ Kč. Začátek \textbf{splátkového kalendáře}:
  \begin{center}
    \renewcommand{\arraystretch}{1.3}
    \begin{tabular}{|c|r|r|r|r|}
      \hline
      rok & dluh na začátku & úrok ($6\,\%$) & úmor & dluh na konci \\
      \hline
      $1$ & $200\,000{,}00$ & $12\,000{,}00$ & $35\,479{,}28$ & $164\,520{,}72$ \\
      $2$ & $164\,520{,}72$ & $9\,871{,}24$ & $37\,608{,}04$ & $126\,912{,}68$ \\
      \vdots & \vdots & \vdots & \vdots & \vdots \\
      \hline
    \end{tabular}
  \end{center}
  Na začátku tvoří velkou část splátky úrok, ke konci skoro jen úmor. U třicetileté hypotéky to znamená, že prvních deset let splácíte hlavně bance, ne dluh.

  \begin{center}
    \begin{tikzpicture}[scale=0.8]
      \begin{scope}
        \draw[->] (0,0) -- (6.3,0) node[right] {\small rok};
        \draw[->] (0,0) -- (0,5.6) node[above] {\small tis. Kč};
        \foreach \y in {1,2,3,4} \draw (0.07,\y) -- (-0.07,\y) node[left] {\tiny $\y0$};
        \foreach \k/\u in {1/1.2, 2/0.987124, 3/0.761476, 4/0.522289, 5/0.268751} {
          \filldraw[fill=gray!40, draw=black] ({\k-0.35},0) rectangle ({\k+0.35},{4.747928-\u});
          \filldraw[fill=colorofwisdom!70, draw=black] ({\k-0.35},{4.747928-\u}) rectangle ({\k+0.35},4.747928);
          \node[below, font=\scriptsize] at (\k,0) {$\k$};
        }
        \filldraw[fill=colorofwisdom!70, draw=black] (0.3,5.15) rectangle (0.6,5.35);
        \node[right, font=\scriptsize] at (0.6,5.25) {úrok};
        \filldraw[fill=gray!40, draw=black] (2.4,5.15) rectangle (2.7,5.35);
        \node[right, font=\scriptsize] at (2.7,5.25) {úmor};
        \node[font=\small] at (3,-1) {složení splátky};
      \end{scope}
      \begin{scope}[xshift=8.3cm]
        \draw[->] (0,0) -- (6.2,0) node[right] {\small rok};
        \draw[->] (0,0) -- (0,5.6) node[above] {\small zbývající dluh};
        \foreach \y/\l in {1/25, 2/50, 3/75, 4/100} \draw (0.07,\y) -- (-0.07,\y) node[left] {\tiny $\l\,\%$};
        \foreach \t in {0,10,20,30} \draw (0.18*\t,0.07) -- (0.18*\t,-0.07) node[below] {\scriptsize $\t$};
        \draw[dashed, gray] (0,4) -- (5.4,0);
        \draw[thick, colorofwisdom, domain=0:30, samples=60] plot ({0.18*\x}, {4*(4.321942-pow(1.05,\x))/3.321942});
        \draw[dashed] (1.8,0) -- (1.8,3.243) -- (0,3.243);
        \filldraw[colorofwisdom] (1.8,3.243) circle (2pt) node[above right, font=\scriptsize] {po $10$ letech zbývá $81\,\%$};
        \node[font=\small] at (2.8,-1) {hypotéka na $30$ let při $5\,\%$};
      \end{scope}
    \end{tikzpicture}
  \end{center}
  Levý graf: všechny splátky jsou stejně vysoké, ale červená část (úrok) se zmenšuje a šedá (úmor) roste. Dluh proto neklesá rovnoměrně, ale zpočátku pomaleji a ke konci rychleji. U pětiletého úvěru je to skoro neznatelné, u dlouhé hypotéky už velmi. Pravý graf ukazuje zbývající dluh u hypotéky na $30$ let při $5\,\%$ ročně: po třetině doby splácení zbývá splatit ještě $81\,\%$ dluhu (dluh po $k$ letech je $U\cdot\frac{1{,}05^{30} - 1{,}05^k}{1{,}05^{30} - 1}$, což plyne ze stejného výpočtu jako v důkazu výše). Šedá přímka ukazuje, jak by dluh klesal, kdyby se splácel rovnoměrně.
\end{example}

\begin{note}\label{note:rpsn}
  Při porovnávání úvěrů se nedívejte jen na úrokovou sazbu, ale na \textbf{RPSN} (roční procentní sazba nákladů). Ta zahrnuje i poplatky za vedení úvěru, pojištění a další náklady. Půjčka `bez úroku' může mít RPSN klidně desítky procent. Nebankovní půjčky v minulosti nabízely RPSN v řádu stovek i tisíců procent, než zákon o spotřebitelském úvěru v roce 2016 výrazně zpřísnil pravidla pro jejich poskytovatele.
\end{note}

\begin{problem}
  \begin{enumerate}[label=(\alph*)]
    \item Telefon za $24\,000$ Kč koupíme na $12$ měsíčních splátek s měsíčním úrokem $1\,\%$. Jaká je splátka a kolik celkem přeplatíme?
    \item Kolik musíme spořit na konci každého roku při $4\,\%$, abychom za $10$ let měli $500\,000$ Kč?
    \item Úvěr $1\,000\,000$ Kč na $20$ let při $5\,\%$ ročně. Jaká je roční splátka a kolik zaplatíme celkem?
  \end{enumerate}
\end{problem}

\begin{solution}
  \begin{enumerate}[label=(\alph*)]
    \item $a = 24000\cdot\frac{0{,}01\cdot 1{,}01^{12}}{1{,}01^{12} - 1}\approx 2\,132{,}37$ Kč. Celkem $25\,588$ Kč, přeplatíme asi $1\,588$ Kč.
    \item $a = 500000\cdot\frac{0{,}04}{1{,}04^{10} - 1}\approx 41\,645$ Kč ročně.
    \item $a = 10^6\cdot\frac{0{,}05\cdot 1{,}05^{20}}{1{,}05^{20} - 1}\approx 80\,243$ Kč ročně, celkem asi $1\,604\,852$ Kč. Na úrocích zaplatíme přes $600\,000$ Kč.
  \end{enumerate}
\end{solution}

\section{Příběh --- Lichváři v sedmém kruhu}

Půjčovat peníze na úrok bylo po většinu evropských dějin hříchem. Už První nikajský koncil v roce 325 zakázal brát úroky duchovním a Třetí lateránský koncil v roce 1179 rozhodl, že lichváři budou exkomunikováni a nesmí být pohřbeni v posvěcené půdě. Za \textit{lichvu} se přitom považoval jakýkoli úrok, ne jen přemrštěný.

Dante Alighieri ve své \textit{Božské komedii} umístil lichváře do sedmého kruhu pekla, mezi ty, kdo se provinili násilím proti Bohu a přírodě. Sedí tam na rozžhaveném písku pod deštěm ohně a každý má na krku měšec se znakem svého rodu, na který neustále zírá. Dante je tam řadí vedle rouhačů a sodomitů. Podle středověké logiky totiž lichvář nechává peníze `rodit' další peníze, což je proti přirozenosti, protože peníze jsou neplodné.

Shakespearův \textit{Kupec benátský} pak vytvořil postavu Shylocka, lichváře, který si jako zástavu za půjčku vymíní libru masa z dlužníkova těla. Dnešní banky naštěstí berou jen nemovitosti. Ale principiálně se toho zas tolik nezměnilo: kdo neumí počítat složené úročení, ten platí. Možná proto je ten citát na začátku kapitoly tak oblíbený, i když ho Einstein nejspíš nikdy neřekl.

\section{Shrnutí a otázky k ústní zkoušce}

\begin{problem}
  Odpovězte vlastními slovy:
  \begin{enumerate}
    \item Co je procento a procentní bod? Proč zdražení o $20\,\%$ a následné zlevnění o $20\,\%$ nevrátí původní cenu?
    \item Jaký je rozdíl mezi jednoduchým a složeným úročením? Jak souvisí s aritmetickou a geometrickou posloupností?
    \item Jak se do výpočtu promítne daň z úroků a častější připisování úroků? Co je efektivní sazba?
    \item Co je inflace a reálný výnos?
    \item Odvoďte vzorec pro pravidelné spoření.
    \item Co je anuitní splácení? Odvoďte vzorec pro výši splátky.
    \item Z čeho se skládá splátka úvěru a jak se poměr mezi úrokem a úmorem v čase mění?
    \item Co je RPSN a proč je důležitější než samotná úroková sazba?
  \end{enumerate}
\end{problem}
